%% ACM sigconf — author draft
\documentclass[sigconf,authordraft]{acmart}
\usepackage{multirow}

\AtBeginDocument{%
  \providecommand\BibTeX{{Bib\TeX}}}

%% Rights / conference placeholder (update before submission)
\setcopyright{acmlicensed}
\copyrightyear{2026}
\acmYear{2026}
\acmDOI{XXXXXXX.XXXXXXX}
\acmConference[NetVulnPy '26]{Seminar on Software Security Analysis}{2026}{University}
\acmISBN{978-1-4503-XXXX-X/2026/01}

%%─────────────────────────────────────────────────────────────────────
\begin{document}

\title{Large-Scale Vulnerability Analysis of Open-Source Python Networking Repositories}

\author{Ahmed Salah}
\affiliation{%
  \institution{University of Passau}
  \city{Passau}
  \country{Germany}}
\email{ahmedelshaikh51@gmail.com}

\renewcommand{\shortauthors}{Ahmed}

%%─── Abstract ────────────────────────────────────────────────────────
\begin{abstract}
Open-source Python repositories focused on networking are widely adopted
in production systems, yet their security posture remains understudied at
scale. This paper presents \textsc{NetVulnPy}, an automated pipeline that
harvests Python networking repositories from GitHub, applies static
application security testing (SAST) using Bandit and Semgrep, and
additionally employs LLM-based analysis for complementary vulnerability
detection. We analyze 1{,}984 repositories using both traditional SAST
tools and an LLM (Qwen3-Next-80B) and report on the distribution of
vulnerabilities by severity, category, and repository characteristics
such as popularity and size. Our results show that high-severity
issues---most prominently insecure subprocess invocation---persist even
in highly starred projects. The LLM analyzer produces 82{,}741 findings
with 46.2\% at high severity, surfacing vulnerability types such as SSRF
and path traversal that require data-flow reasoning beyond pattern
matching. Because raw Bandit output is heavily dominated by a single
low-severity rule (\texttt{assert\_used}, 86.9\% of all findings), we
report both raw and noise-excluded counts and treat all detections as
\emph{potential} rather than confirmed vulnerabilities. The three tools
exhibit markedly different rule profiles: Bandit excels at
Python-idiomatic security smells, Semgrep provides targeted taint-tracking
coverage, and the LLM contributes semantic reasoning at the cost of
taxonomic consistency. The pipeline and all data are released as
open-source artifacts to support reproducibility.
\end{abstract}

%%─── CCS Concepts (generate yours at https://dl.acm.org/ccs/ccs.cfm) ─
\begin{CCSXML}
<ccs2012>
 <concept>
  <concept_id>10002978.10003029.10003032</concept_id>
  <concept_desc>Security and privacy~Software security engineering</concept_desc>
  <concept_significance>500</concept_significance>
 </concept>
 <concept>
  <concept_id>10011007.10011074.10011092</concept_id>
  <concept_desc>Software and its engineering~Software verification and validation</concept_desc>
  <concept_significance>300</concept_significance>
 </concept>
</ccs2012>
\end{CCSXML}
\ccsdesc[500]{Security and privacy~Software security engineering}
\ccsdesc[300]{Software and its engineering~Software verification and validation}

\keywords{static analysis, vulnerability detection, Python, open source,
  Bandit, Semgrep, LLM, SAST, software security}

\maketitle

%%═════════════════════════════════════════════════════════════════════
\section{Introduction}
\label{sec:intro}

Python has become one of the most widely used programming languages,
ranking consistently in the top three of major language indices. Its
rich ecosystem of networking libraries---including \texttt{requests},
\texttt{aiohttp}, \texttt{paramiko}, and \texttt{scapy}---powers
critical infrastructure ranging from web services to network
monitoring tools. However, the very openness that accelerates
development also introduces security risks: developers frequently
reuse code with insufficient scrutiny, and known vulnerability
patterns propagate across repositories. The scale of this problem is
substantial: Alfadel et al.~\cite{Alfadel2023} found that security
vulnerabilities in the PyPI ecosystem take a median of 39~months to
be discovered, and that basic networking and utility packages are
among the highest risk factors for downstream projects---with the
HTTP client package \texttt{urllib3} alone exposing 23.83\% of
analyzed Python projects to known vulnerabilities. Moreover, PyPI
currently allows developers to publish packages without any mandatory
security checks or code vetting~\cite{Alfadel2023}, highlighting the
need for proactive, automated vulnerability detection.

Static Application Security Testing (SAST) tools offer a scalable
means of identifying vulnerabilities without executing code. Tools
such as Bandit and Semgrep can detect common weakness
patterns---from SQL injection and command injection to hardcoded
credentials and insecure cryptographic usage---by analyzing source
code structure and data flow. Despite the availability of these
tools, prior work has shown that many open-source projects do not
integrate SAST into their development
workflows~\cite{Kaur2020, Ruohonen2021}.

This paper addresses the following research questions:

\begin{description}
  \item[\textbf{RQ1.}] What types and severities of vulnerabilities
    are most prevalent in open-source Python networking repositories?
  \item[\textbf{RQ2.}] Is there a relationship between repository
    characteristics (stars, size, age, activity) and the density of
    detected vulnerabilities?
  \item[\textbf{RQ3.}] How do Bandit and Semgrep compare in terms of
    detection coverage and complementarity on the same codebase?
\end{description}

To answer these questions, we designed and implemented
\textsc{NetVulnPy}, a four-stage automated pipeline that:
(1)~harvests Python networking repositories from the GitHub Search
API, (2)~downloads and extracts their source files,
(3)~runs Bandit and Semgrep on each repository, and
(4)~loads all findings into a SQLite database for structured
querying and visualization.

\textbf{Contributions.} Relative to prior single-tool, PyPI-snapshot,
or dependency-chain studies, this paper makes three new empirical
contributions: (i)~it is, to our knowledge, the first study to focus
SAST analysis specifically on the Python \emph{networking} domain at
this scale; (ii)~it analyzes \emph{first-party} source code in repositories'
active development state on GitHub rather than released PyPI snapshots or
mined third-party vulnerability reports; and (iii)~it applies two
complementary SAST tools (Bandit and Semgrep) to an identical corpus,
enabling a direct comparison of their rule profiles and severity
distributions. We also confirm, on this new corpus, two patterns previously
reported on PyPI---the dominance of low-severity \texttt{assert} noise and
the weak relationship between popularity and finding density.

The remainder of this paper is organized as follows.
Section~\ref{sec:background} reviews related work on vulnerability
analysis and SAST tools. Section~\ref{sec:methodology} describes the
\textsc{NetVulnPy} pipeline and experimental setup.
Section~\ref{sec:results} presents our findings.
Section~\ref{sec:discussion} discusses implications, threats to
validity, and tool comparison. Section~\ref{sec:conclusion} concludes
with future directions.


%%═════════════════════════════════════════════════════════════════════
\section{Background and Related Work}
\label{sec:background}

\subsection{Vulnerabilities in Open-Source Python Software}

The prevalence of security vulnerabilities in open-source Python
software has attracted growing research attention. Ruohonen~\cite{Ruohonen2018}
conducted an early empirical analysis of vulnerabilities in Python
packages for web applications, establishing that security issues are
widespread even in popular packages. Building on this,
Alfadel et al.~\cite{Alfadel2023} performed a large-scale empirical
analysis of security vulnerabilities across Python packages, examining
how vulnerabilities are introduced, persist, and eventually get fixed.
Analyzing 2,224 open-source Python GitHub repositories and mining
vulnerability reports from the Snyk.io database, they found that
the majority of Python vulnerabilities are of medium severity
(64.03\%), followed by high severity (25.42\%), with Cross-Site
Scripting (XSS), Denial of Service (DoS), and Information Exposure
being the most common types. Notably, their analysis focused on
\emph{dependency-chain} vulnerabilities---security issues introduced
through third-party packages that a project imports---rather than
\emph{first-party} vulnerabilities present in the project's own
source code. They also found that basic networking packages such as
\texttt{urllib3} are among the biggest offenders, exposing 23.83\% of
analyzed projects to vulnerabilities. The authors concluded that
current tools are often not qualified to find complex issues and
called for ``robust tools that combine exhaustive techniques like
program analysis, testing, and verification.'' Our work complements
theirs by using Bandit and Semgrep to detect first-party source-code
vulnerabilities directly, and we compare our severity distributions
against theirs as a baseline.

Most closely related to our work, Ruohonen et al.~\cite{Ruohonen2021}
performed a large-scale security-oriented static analysis of over
197,000 Python packages in PyPI using Bandit, detecting more than
749,000 security issues across 46\% of packages. They found that
generic exception handling (\texttt{try-except-pass}) and command
injection via the \texttt{subprocess} module were the most prevalent
vulnerability types, with low-severity issues dominating but medium
and high severity issues accounting for approximately 41\% of
findings. Notably, they disabled Bandit's assertion detector (B101)
due to excessive false positives, a design choice we consider in our
own methodology. The authors identified three key limitations of their
study: (1)~reliance on a single SAST tool, (2)~scanning all of PyPI
without domain-specific filtering, and (3)~analyzing static package
snapshots rather than active development repositories. Our work
directly addresses all three limitations by employing both Bandit and
Semgrep, focusing specifically on networking-domain repositories, and
analyzing projects sourced from GitHub in their active development
state.

More recently, Quan et al.~\cite{Quan2025} presented an empirical
study of vulnerabilities in Python packages and their detection,
confirming that the landscape of Python security issues continues to
evolve and that detection techniques must keep pace with emerging
vulnerability patterns.

The OWASP Top Ten~\cite{OWASP2021} provides a widely adopted
classification of the most critical web application security risks,
including injection, broken authentication, and security
misconfiguration. Many of these categories are detectable through
static analysis of Python source code.

\subsection{Static Application Security Testing (SAST)}

SAST tools analyze source code without execution to identify
potential security flaws. Kaur and Nayyar~\cite{Kaur2020} conducted a
comparative study of static code analysis tools for vulnerability
detection in C/C++ and Java, demonstrating that different tools
exhibit varying strengths across vulnerability categories and that no
single tool provides complete coverage. Their findings motivate the
multi-tool approach adopted in our work.

\textbf{Bandit} is a Python-specific SAST tool maintained by the
PyCQA community. It applies a set of predefined plugins (test IDs
such as \texttt{B101}--\texttt{B703}) that pattern-match against known
insecure Python constructs. Bandit groups its detectors into seven
categories, including one specifically targeting insecure network
protocols (TLS misconfiguration, certificate validation bypass,
insecure SSH host keys, and clear-text protocols such as FTP and
Telnet)~\cite{Ruohonen2021}. Each finding is annotated with both a
severity level and a \emph{confidence} level, the latter of which can
be used to suppress noisy detectors. Bandit is lightweight, fast, and
produces structured JSON output, making it suitable for automated
pipelines.

\textbf{Semgrep} is a multi-language static analysis tool that uses
pattern-matching with a syntax-aware engine. Unlike Bandit, Semgrep
supports user-defined rules and inter-procedural taint tracking. Its
\texttt{p/python} ruleset includes security, correctness, and
performance checks drawn from community contributions.

As Kaur and Nayyar~\cite{Kaur2020} demonstrated, no single tool
achieves complete coverage across all vulnerability categories. We
apply this principle by combining Bandit and Semgrep to maximize
detection breadth.

\subsection{Mining GitHub for Security Research}

GitHub is the largest host of open-source code and a common data
source for empirical software engineering research. Our methodology
addresses known biases in repository sampling---such as the
prevalence of inactive, toy, or forked projects---by filtering for
Python-language repositories with networking-related keywords, sorting
by popularity, and reporting repository characteristics alongside
vulnerability counts. Unlike prior work that analyzed static PyPI
snapshots~\cite{Ruohonen2021}, our GitHub-based approach captures
projects in their active development environment, including unreleased
code and work in progress.


%%═════════════════════════════════════════════════════════════════════
\section{Methodology}
\label{sec:methodology}

\subsection{Pipeline Overview}

\textsc{NetVulnPy} is an end-to-end Python pipeline that
(1)~harvests Python networking repositories from the GitHub Search
API, (2)~downloads and extracts their Python source files,
(3)~runs Bandit, Semgrep, and/or an LLM-based analyzer on each
repository, and (4)~loads all findings into a SQLite database for
structured querying and visualization via an interactive Streamlit
dashboard.

We query the GitHub Search API for repositories matching
\texttt{language:python} combined with networking-related keywords
(\texttt{socket}, \texttt{http}, \texttt{network}, \texttt{ssl},
\texttt{dns}, \texttt{proxy}), sorted by star count to prioritize
actively maintained projects. For each repository, we record metadata
including star count, fork count, creation date, last update date,
license, and size. Source code is downloaded as ZIP archives via the
GitHub API, and only \texttt{.py} files are extracted.

Each repository is analyzed using the selected tool. Bandit is
invoked as \texttt{bandit -r <repo> -f json -q}; Semgrep is invoked
as \texttt{semgrep scan --config p/python --no-git-ignore --json
--quiet <repo>}. Both produce JSON reports containing individual
findings with severity, affected file and line number, and a code
snippet. Semgrep severity levels (ERROR, WARNING, INFO) are mapped to
(HIGH, MEDIUM, LOW) for consistency with Bandit. Repositories with no
\texttt{.py} files are skipped, and previously analyzed repositories
are not re-analyzed, enabling incremental runs.

The LLM-based analyzer sends each Python file (up to 50 files per
repository, truncated at 15{,}000 characters) to a large language
model via an OpenAI-compatible API. The model is prompted to return
a structured JSON array of findings following the same schema as the
SAST tools. We used \texttt{Qwen3-Next-80B-A3B-Instruct-FP8}
with temperature set to 0.1 to minimize non-determinism. The LLM
self-assigns severity, confidence, and generates descriptive rule
identifiers for each finding.

\subsection{Experimental Setup}

We harvested the top 2{,}000 Python networking repositories by star
count from GitHub (as of May 2026). After filtering repositories with
no Python files, 1{,}984 repositories remained for analysis. Each
repository was analyzed independently with both Bandit (v1.7+) and
Semgrep (v1.99+). The pipeline achieved a 100\% scan success rate
across all repositories for both tools. The pipeline was executed on
a macOS system with Python~3.13 using 8 parallel workers.


Lines of code (LOC) per repository, used in the vulnerability-density
metric (Section~\ref{sec:results}), are counted as \emph{logical}
lines: all non-blank, non-comment lines across every extracted
\texttt{.py} file, where a comment line is defined as one whose
stripped content begins with \texttt{\#}. This counting is performed
by a custom function in the pipeline (no external tool used)

All source code, the analysis pipeline, and the resulting database
are available as
\href{https://github.com/ahmedelshaikh20/NetVulnPy}{open-source artifacts}.
%%═════════════════════════════════════════════════════════════════════
\section{Results}
\label{sec:results}

Bandit was applied to all 1{,}984 repositories with a 100\% scan
success rate, yielding a total of 698{,}020 findings across 72
distinct Bandit test IDs. Of the scanned repositories, 35.2\%
contained at least one high-severity finding. As discussed below, the
great majority of these findings stem from a single low-severity rule;
we therefore report both raw and noise-excluded counts throughout.

\subsection{Overall Vulnerability Landscape (RQ1)}

Table~\ref{tab:severity-bandit} presents the distribution of findings
by severity level.

\begin{table}[h]
  \caption{Vulnerability counts by severity (Bandit).}
  \label{tab:severity-bandit}
  \begin{tabular}{lrr}
    \toprule
    Severity & Count & \% of Total \\
    \midrule
    HIGH   &   4{,}866 &  0.7\% \\
    MEDIUM &  24{,}131 &  3.5\% \\
    LOW    & 669{,}023 & 95.8\% \\
    \midrule
    Total  & 698{,}020 & 100\% \\
    \bottomrule
  \end{tabular}
\end{table}

Low-severity issues overwhelmingly dominate the dataset at 95.8\%,
driven primarily by the \texttt{assert\_used} rule (B101), which alone
accounts for 606{,}297 findings (86.9\% of all detections). Excluding
B101, the corpus contains 91{,}723 Bandit findings, and we report this
B101-excluded figure alongside the raw total wherever it materially
affects interpretation. This pattern is consistent with findings from
Ruohonen et al.~\cite{Ruohonen2021}, who also observed that
low-severity issues dominate Bandit output, leading them to disable
B101 due to excessive noise. Despite the low-severity dominance, the
absolute count of high-severity findings (4{,}866) across 35.2\% of
repositories indicates that serious potential security issues remain
widespread in networking code.

\textbf{Most common vulnerability categories.}
Table~\ref{tab:top-rules} lists the ten most frequently triggered
Bandit rules across all repositories.

\begin{table}[h]
  \caption{Top-10 most frequent Bandit rules.}
  \label{tab:top-rules}
  \begin{tabular}{rllrr}
    \toprule
    Rank & ID & Rule Name & Count & Sev. \\
    \midrule
    1  & B101 & assert\_used                        & 606{,}297 & LOW \\
    2  & B105 & hardcoded\_password\_string          &  13{,}502 & LOW \\
    3  & B311 & blacklist (random)                   &  12{,}429 & LOW \\
    4  & B110 & try\_except\_pass                    &  10{,}202 & LOW \\
    5  & B603 & subprocess\_without\_shell\_true     &   8{,}044 & LOW \\
    6  & B607 & start\_process\_with\_partial\_path  &   4{,}794 & LOW \\
    7  & B404 & blacklist (import subprocess)        &   4{,}095 & LOW \\
    8  & B113 & request\_without\_timeout            &   4{,}048 & MED \\
    9  & B106 & hardcoded\_password\_funcarg         &   3{,}978 & LOW \\
    10 & B108 & hardcoded\_tmp\_directory            &   3{,}125 & MED \\
    \bottomrule
  \end{tabular}
\end{table}

The most prevalent finding, B101 (\texttt{assert\_used}), flags the use
of \texttt{assert} statements which are stripped in optimized bytecode,
potentially disabling security checks at runtime. While individually
low-risk, the sheer volume (606{,}297 instances) indicates a
pervasive idiom rather than a security crisis. Beyond B101,
subprocess-related rules (B603, B607, B404) collectively account for
16{,}933 findings. These rules flag subprocess usage that \emph{may}
enable command injection (OWASP A03:2021~\cite{OWASP2021}); we note,
however, that B603 (the largest of the three) specifically detects
subprocess calls made \emph{without} \texttt{shell=True}---the
comparatively safer invocation form---and is rated LOW severity. The
higher-risk \texttt{shell=True} pattern is captured by a separate rule
(B602), which does not appear among the most frequent rules in our
corpus. These subprocess counts should therefore be read as
\emph{exposure surface} rather than confirmed injection vectors.
Rules B105 and B106 flag 17{,}480 \emph{candidate} hardcoded secrets;
both are LOW-severity, pattern-based detectors that match string
literals near password-like identifiers and are known to produce many
false positives (test fixtures, example configuration, placeholder
values), so without manual triage these counts indicate potential
rather than confirmed credential exposure. The presence of
\texttt{request\_without\_timeout} (B113) is particularly relevant to
the networking domain, as missing timeouts can lead to denial-of-service
conditions in client applications.


\subsection{Repository Characteristics and Vulnerability Density (RQ2)}

To investigate whether repository popularity or size correlates with
vulnerability density, we compute the metric:
\begin{equation}
  \text{vuln\_density} = \frac{\text{total\_issues}}{\text{LOC}} \times 1000
  \label{eq:density}
\end{equation}
expressing vulnerabilities per thousand lines of code (findings/KLOC).
We caution that, for the Bandit-derived density, \texttt{total\_issues}
is dominated by the low-severity B101 rule (86.9\% of all Bandit
findings); the metric therefore partly reflects \texttt{assert} usage
rather than security-relevant patterns.
%% TODO: recompute Equation~\ref{eq:density} and the correlations in
%% Table~\ref{tab:correlations} with B101 EXCLUDED, and ideally also for
%% HIGH-severity findings only, to isolate security-relevant density.
%% Report whether the correlation signs/magnitudes change. The current
%% coefficients below are computed on the B101-inclusive total.

The distribution of vulnerability density across repositories is
heavily right-skewed: the majority of repositories exhibit fewer than
50~findings per KLOC, while a long tail extends beyond
300~findings per KLOC, indicating a small number of repositories with
exceptionally high vulnerability density.

\textbf{Popularity vs.\ Vulnerability Density.}
We computed Pearson and Spearman correlation coefficients between
repository star count and vulnerability density
(Table~\ref{tab:correlations}). The Pearson correlation is
$r = 0.208$  and the Spearman rank correlation is
$\rho = 0.373$ . Both coefficients indicate a weak
\emph{positive} correlation: more popular repositories tend to exhibit
\emph{marginally higher} vulnerability density, not lower. The effect
size is small, popularity explains only about 4\% of the variance in
density ($r^2 \approx 0.04$)---so we interpret this as evidence that
repository popularity is \emph{not} a reliable indicator of better code
quality, and if anything the relationship runs weakly in the opposite
direction.
% TODO: VERIFY THE SIGN of this correlation against the raw data. The
% earlier draft described a scatter plot in which the highest-density
% repositories were small/low-star, which would imply a NEGATIVE
% relationship and is INCONSISTENT with the positive coefficients
% reported here. Either the coefficients or that scatter description
% was wrong; resolve before submission and fix the prose accordingly.

\textbf{Repository Size vs.\ Vulnerability Density.}
The correlation between repository size (in KB) and vulnerability
density is negligible: Pearson $r = -0.015$ ($p = 0.502$, not
statistically significant), Spearman $\rho = 0.050$ ($p = 0.026$).
Larger repositories do not exhibit systematically higher or lower
vulnerability density. With $n = 1{,}984$, even trivially small effects
can reach statistical significance; the Spearman result here
($\rho = 0.050$, $p = 0.026$) is significant but practically null, and
we do not interpret it as a meaningful relationship. This is consistent
with the expectation that vulnerability density normalizes for project
size.

\textbf{Repository Age vs.\ Findings.}
Analysis of repository age (in days since creation) against total
findings reveals no clear temporal trend. Findings are distributed
across repositories of all ages, with high-finding outliers appearing
at every age bracket. Older repositories do not necessarily accumulate
more findings than newer ones, suggesting that vulnerability
introduction is not a function of project maturity.

\begin{table}[h]
  \caption{Correlation between repository metadata and vulnerability density (findings/KLOC, B101-inclusive).}
  \label{tab:correlations}
  \begin{tabular}{lrrrr}
    \toprule
    Variable & Pearson $r$ & $p$-value & Spearman $\rho$ & $p$-value \\
    \midrule
    Stars (popularity) & 0.208 & $<$0.001 & 0.373 & $<$0.001 \\
    Size (KB)          & $-$0.015 & 0.502 & 0.050 & 0.026 \\
    \bottomrule
  \end{tabular}
\end{table}

\newpage
\subsection{Semgrep Results}

To complement the Bandit analysis, we ran Semgrep with the
\texttt{p/python} community ruleset on the same 1{,}984 repositories.
Table~\ref{tab:severity-semgrep} summarizes the severity distribution.

\begin{table}[h]
  \caption{Vulnerability counts by severity (Semgrep).}
  \label{tab:severity-semgrep}
  \begin{tabular}{lrr}
    \toprule
    Severity & Count & \% of Total \\
    \midrule
    HIGH   & 1{,}678 & 29.9\% \\
    MEDIUM & 3{,}427 & 61.0\% \\
    LOW    &    509  &  9.1\% \\
    \midrule
    Total  & 5{,}614 & 100\% \\
    \bottomrule
  \end{tabular}
\end{table}

Semgrep detected 5{,}614 findings across 79 distinct rules, with
19.7\% of repositories containing at least one high-severity finding.
Unlike Bandit, where low-severity issues dominate due to B101
(\texttt{assert\_used}), Semgrep's severity distribution is
concentrated at the medium level (61.0\%), followed by high (29.9\%).

Table~\ref{tab:top-rules-semgrep} lists the ten most frequently
triggered Semgrep rules.

\begin{table}[h]
  \caption{Top-10 most frequent Semgrep rules.}
  \label{tab:top-rules-semgrep}
  \begin{tabular}{rlr}
    \toprule
    Rank & Rule & Sev. \\
    \midrule
    1  & subprocess-shell-true              & HIGH \\
    2  & insecure-hash-algorithm-md5        & MED \\
    3  & logger-credential-disclosure       & MED \\
    4  & insecure-hash-algorithm-sha1       & MED \\
    5  & request-with-http                  & MED \\
    6  & insecure-file-permissions          & MED \\
    7  & dangerous-subprocess-tainted-env   & HIGH \\
    8  & insecure-uuid-version              & LOW \\
    9  & bind-to-all-interfaces             & MED \\
    10 & ssl-wrap-socket-deprecated         & MED \\
    \bottomrule
  \end{tabular}
\end{table}

\textbf{Correlation analysis.}
Semgrep findings show even weaker correlations with repository
metadata than Bandit: stars vs.\ vulnerability density yields
Pearson $r = 0.134$ ($p < 0.001$), Spearman $\rho = 0.160$
($p < 0.001$); repository size vs.\ density yields Pearson
$r = -0.003$ ($p = 0.880$), Spearman $\rho = 0.109$ ($p < 0.001$).
These results confirm that neither popularity nor size strongly
predicts vulnerability density regardless of the tool used.

\subsection{LLM-Based Analysis}

To explore whether large language models (LLMs) can serve as a
complementary vulnerability detection mechanism, we extended the
\textsc{NetVulnPy} pipeline with an LLM-based analyzer. Each Python
file in a repository is sent to an LLM with a structured prompt
requesting security analysis; the model returns findings in a
standardized JSON format identical to the SAST tool output schema. We
used \texttt{Qwen3-Next-80B-A3B-Instruct} (8-bit quantized), hosted on
a university inference endpoint compatible with the OpenAI API.

The LLM analyzer successfully processed 1{,}945 of 2{,}000 repositories
(97.3\% success rate; 55 failures due to API errors or unparseable
responses), producing 82{,}741 findings across 14{,}693 distinct rule
identifiers. Table~\ref{tab:severity-llm} shows the severity
distribution.

\begin{table}[h]
  \caption{Vulnerability counts by severity (LLM).}
  \label{tab:severity-llm}
  \begin{tabular}{lrr}
    \toprule
    Severity & Count & \% of Total \\
    \midrule
    HIGH   & 38{,}231 & 46.2\% \\
    MEDIUM & 37{,}044 & 44.8\% \\
    LOW    &  7{,}466 &  9.0\% \\
    \midrule
    Total  & 82{,}741 & 100\% \\
    \bottomrule
  \end{tabular}
\end{table}

The LLM's severity distribution is strikingly different from both
traditional SAST tools: high-severity findings constitute 46.2\% of all
detections (vs.\ 0.7\% for Bandit and 29.9\% for Semgrep), with 95.9\%
of repositories flagged as containing at least one high-severity issue.
This suggests the LLM applies a more aggressive severity classification,
potentially reflecting its training on security literature where
vulnerability descriptions tend to emphasize worst-case impact.

\textbf{Most common LLM-detected vulnerability categories.}
Table~\ref{tab:top-rules-llm} lists the ten most frequently flagged
categories.

\begin{table}[h]
  \caption{Top-10 most frequent LLM-detected vulnerability categories.}
  \label{tab:top-rules-llm}
  \begin{tabular}{rlrr}
    \toprule
    Rank & Category & Count & Sev. \\
    \midrule
    1  & Hardcoded Secrets            & 8{,}477 & HIGH \\
    2  & Unsanitized Network Input    & 3{,}738 & MED  \\
    3  & Path Traversal               & 2{,}793 & HIGH \\
    4  & Injection (JSON/Input)       & 2{,}544 & HIGH \\
    5  & Hardcoded File Paths         & 1{,}834 & HIGH \\
    6  & Command Injection (Env Vars) & 1{,}756 & HIGH \\
    7  & SQL Injection                & 1{,}559 & HIGH \\
    8  & Command Injection (Shell)    & 1{,}486 & HIGH \\
    9  & SSRF                         & 1{,}201 & HIGH \\
    10 & Code Injection               & 1{,}151 & HIGH \\
    \bottomrule
  \end{tabular}
\end{table}

The LLM's top findings are dominated by high-severity categories. The
most frequent category, hardcoded secrets (8{,}477 findings), overlaps
conceptually with Bandit's B105/B106 rules (17{,}480 combined findings),
but the LLM flags fewer instances, potentially because it can use
semantic context to distinguish real credentials from test fixtures and
placeholder values---a distinction that pattern-based rules cannot make.
Command injection categories (ranks 6, 8, 10) collectively account for
4{,}378 findings, compared to 16{,}933 subprocess-related Bandit
findings; the lower count may reflect the LLM's ability to assess
whether user-controlled input actually reaches the subprocess call. The
LLM also surfaces vulnerability types not covered by the default Bandit
or Semgrep rulesets, notably SSRF (1{,}201 findings) and path traversal
(2{,}793 findings), which require data-flow reasoning beyond simple
pattern matching.

\textbf{Correlation analysis.}
The LLM findings exhibit a weak \emph{negative} correlation between
repository popularity and vulnerability density: Pearson $r = -0.139$
($p < 0.001$), Spearman $\rho = -0.385$ ($p < 0.001$). This is the
opposite sign from the Bandit and Semgrep correlations (both weakly
positive), suggesting that the LLM may be more sensitive to code quality
patterns that improve with community scrutiny. Repository size vs.\
density remains negligible: Pearson $r = -0.015$ ($p = 0.516$), Spearman
$\rho = -0.435$ ($p < 0.001$).

\textbf{Limitations of LLM-based analysis.}
While the LLM generates a rich and diverse set of findings, several
caveats apply: (1)~the 14{,}693 distinct ``rule'' identifiers are
LLM-generated labels, not standardized taxonomies, making cross-tool
comparison imprecise; (2)~LLMs are known to hallucinate, so some
findings may describe vulnerabilities that do not exist in the code;
(3)~the analysis is limited by context window size (files are truncated
at 15{,}000 characters) and per-file granularity, preventing
cross-file taint analysis; (4)~reproducibility is limited by model
non-determinism (temperature was set to 0.1 to minimize variance). We
therefore treat LLM findings as a \emph{complementary signal} requiring
validation rather than a replacement for established SAST tools.

\subsection{Tool Comparison: Bandit vs.\ Semgrep vs.\ LLM (RQ3)}

Table~\ref{tab:tool-comparison} contrasts the three tools on the same
corpus.

\begin{table}[h]
  \caption{Detection comparison: Bandit vs.\ Semgrep vs.\ LLM.}
  \label{tab:tool-comparison}
  \begin{tabular}{lrrr}
    \toprule
    Metric & Bandit & Semgrep & LLM \\
    \midrule
    Total findings            & 698{,}020 & 5{,}614 & 82{,}741 \\
    Total (excl.\ B101)       & 91{,}723  & 5{,}614 & 82{,}741 \\
    Distinct rules            & 72        & 79      & 14{,}693 \\
    Repos with $\geq$1 HIGH   & 35.2\%   & 19.7\%  & 95.9\% \\
    HIGH findings             & 4{,}866   & 1{,}678  & 38{,}231 \\
    MEDIUM findings           & 24{,}131  & 3{,}427  & 37{,}044 \\
    LOW findings              & 669{,}023 & 509      & 7{,}466 \\
    Scan success rate         & 100\%     & 100\%    & 97.3\% \\
    \bottomrule
  \end{tabular}
\end{table}

Bandit produces the highest raw finding count, primarily because
Bandit's B101 (\texttt{assert\_used}) rule alone generates 606{,}297
low-severity findings. Excluding B101, Bandit produces 91{,}723
findings, comparable in scale to the LLM's 82{,}741 findings.
Semgrep's community rules are more targeted, producing 5{,}614
findings---an order of magnitude fewer. The LLM occupies a middle
ground in volume but diverges sharply in severity distribution: 46.2\%
of its findings are high-severity, compared to 0.7\% for Bandit and
29.9\% for Semgrep.

The three tools exhibit distinct detection profiles. Bandit excels at
flagging Python-idiomatic security smells (assert usage, subprocess
patterns, hardcoded credentials via pattern matching) but produces
substantial noise. Semgrep provides more targeted detections with
taint-tracking capabilities and a higher proportion of actionable
findings. The LLM contributes a qualitatively different perspective:
it can reason about semantic context (e.g., distinguishing test
fixtures from real credentials), surfaces vulnerability types
requiring data-flow reasoning (SSRF, path traversal), and generates
vulnerability descriptions in natural language. However, the LLM's
14{,}693 distinct ``rule'' identifiers---compared to Bandit's 72 and
Semgrep's 79---reflect a lack of taxonomic consistency, and its 95.9\%
high-severity repository rate suggests over-classification.

We caution that cross-tool severity comparisons are limited by
differing severity semantics: Bandit uses its own calibrated severity
scale, Semgrep's ERROR/WARNING/INFO levels are mapped to
HIGH/MEDIUM/LOW, and the LLM self-assigns severity based on its
training data. The LLM's substantially higher high-severity rate
likely reflects both genuine detection of issues missed by pattern
matching and over-reporting driven by the model's tendency to
emphasize worst-case impact.

\subsection{Precision Validation}

To estimate the true-positive rate of each tool, we performed an
LLM-assisted validation of a stratified random sample of 100~findings
per analyzer. For each sampled finding, we retrieved the full source
file from the downloaded repository and submitted both the tool's
report and the complete source code to an independent LLM
(\texttt{Qwen3-Next-80B-A3B-Instruct}) for verification. The verifier
was prompted to classify each finding as \textsc{true\_positive} (a
genuine, exploitable vulnerability in production context) or
\textsc{false\_positive} (a false alarm---safe usage, test code, or
non-exploitable pattern), along with a confidence level and reasoning.
Each 100-finding sample was stratified across severity levels (34~HIGH,
33~MEDIUM, 33~LOW) to ensure representation at every severity tier.
Pip-audit findings were excluded from this validation because they
flag dependency-level CVEs rather than first-party source code, and
the corresponding source files are not present in the downloaded
repositories.

Table~\ref{tab:precision} summarizes the results.

\begin{table}[h]
  \caption{Precision validation: 100 stratified samples per analyzer.}
  \label{tab:precision}
  \begin{tabular}{lrrr}
    \toprule
    Analyzer & Evaluated & True Pos. & Precision \\
    \midrule
    Semgrep  & 100 & 21 & 21.0\% \\
    Bandit   & 100 & 10 & 10.0\% \\
    LLM      & 100 &  9 &  9.0\% \\
    \bottomrule
  \end{tabular}
\end{table}

Semgrep achieves the highest overall precision at 21.0\%, followed by
Bandit (10.0\%) and the LLM (9.0\%). The differences are more
pronounced at the HIGH severity level
(Table~\ref{tab:precision-severity}): Semgrep confirms 47.1\% of its
HIGH findings as true positives, compared to 29.4\% for Bandit and
20.6\% for the LLM. At MEDIUM severity, precision drops sharply for
all tools (Semgrep: 15.2\%, LLM: 6.1\%, Bandit: 0.0\%). No tool
achieved any true positives among LOW-severity findings in the sample,
confirming that low-severity detections are overwhelmingly noise.

\begin{table}[h]
  \caption{Precision by reported severity level (100 samples per tool).}
  \label{tab:precision-severity}
  \begin{tabular}{llrr}
    \toprule
    Severity & Analyzer & TP / Evaluated & Precision \\
    \midrule
    \multirow{3}{*}{HIGH}
      & Semgrep & 16 / 34 & 47.1\% \\
      & Bandit  & 10 / 34 & 29.4\% \\
      & LLM     &  7 / 34 & 20.6\% \\
    \midrule
    \multirow{3}{*}{MEDIUM}
      & Semgrep &  5 / 33 & 15.2\% \\
      & LLM     &  2 / 33 &  6.1\% \\
      & Bandit  &  0 / 33 &  0.0\% \\
    \midrule
    \multirow{3}{*}{LOW}
      & Semgrep &  0 / 33 &  0.0\% \\
      & Bandit  &  0 / 33 &  0.0\% \\
      & LLM     &  0 / 33 &  0.0\% \\
    \bottomrule
  \end{tabular}
\end{table}

These results have several implications. First, all three tools
exhibit high false-positive rates (79--91\%), confirming that raw SAST
output requires triage before it can inform remediation decisions.
Second, Semgrep's higher precision---particularly at HIGH
severity---validates its design philosophy of fewer, more targeted
rules over broad pattern matching. Third, the LLM's precision (9.0\%)
is comparable to Bandit's (10.0\%) despite producing qualitatively
different findings; the LLM's confirmed true positives include
categories absent from Bandit's ruleset (SQL injection, path
traversal, SSRF), demonstrating complementary value even at modest
precision. Fourth, the uniformly zero precision at LOW severity across
all tools strongly supports filtering or suppressing low-severity
findings in practice.

We note that this validation itself uses an LLM as the verifier, which
introduces potential bias: the verifier may share systematic blind
spots with the LLM analyzer. We discuss this limitation further in
Section~\ref{sec:discussion}.


%%═════════════════════════════════════════════════════════════════════
\section{Discussion}
\label{sec:discussion}

\subsection{Key Findings}

Our analysis of 1{,}984 Python networking repositories reveals several
important observations:

\begin{enumerate}
  \item \textbf{High-severity issues are widespread but mostly
    unconfirmed.} Over one in three repositories (35.2\%) contains at
    least one high-severity Bandit finding, with 4{,}866 high-severity
    detections in total. However, our precision validation
    (Section~\ref{sec:results}) reveals that only 29.4\% of Bandit's
    HIGH findings and 47.1\% of Semgrep's HIGH findings are confirmed
    true positives. Subprocess-related rules (B603, B607, B404)
    collectively flag 16{,}933 subprocess call sites (potential
    injection surface; OWASP A03:2021~\cite{OWASP2021}), though the
    most frequent of these (B603) detects the lower-risk
    non-\texttt{shell} invocation form. These counts therefore
    represent exposure surface rather than confirmed vulnerabilities.

  \item \textbf{Low-severity noise dominates raw counts.} The
    \texttt{assert\_used} rule (B101) alone produces 606{,}297
    findings---86.9\% of all detections. This confirms the observation
    by Ruohonen et al.~\cite{Ruohonen2021} that B101 generates
    excessive noise. Bandit additionally records a \emph{confidence}
    level for every finding (stored in our schema but not yet used for
    filtering); confidence-based filtering offers a principled,
    zero-cost way to suppress B101-class noise and should be applied in
    future analyses.

  \item \textbf{Candidate hardcoded credentials are common.} Rules
    B105 and B106 together flag 17{,}480 candidate hardcoded secrets.
    Because these pattern-based, low-severity rules are prone to false
    positives, this figure should be read as an upper bound on
    potential credential exposure rather than confirmed poor practice;
    manual triage (see Threats to Validity) is needed to estimate the
    true rate.

  \item \textbf{The three tools have distinct detection profiles.}
    Bandit produces 698{,}020 findings (91{,}723 excluding B101 noise),
    Semgrep produces 5{,}614 more targeted findings with 29.9\%
    at high severity, and the LLM produces 82{,}741 findings with
    46.2\% at high severity. The three tools surface largely different
    top rules and vulnerability categories, suggesting that no single
    tool provides complete coverage and supporting the multi-tool
    approach advocated by Kaur and Nayyar~\cite{Kaur2020}. The LLM
    additionally surfaces vulnerability types (SSRF, path traversal)
    that require data-flow reasoning beyond pattern matching.

  \item \textbf{LLM analysis offers complementary but noisy coverage.}
    The LLM generates 14{,}693 distinct rule identifiers (vs.\ 72 for
    Bandit and 79 for Semgrep), reflecting rich but taxonomically
    inconsistent output. Its high-severity rate (95.9\% of repos flagged)
    suggests over-classification, yet the semantic reasoning capability
    allows it to surface categories absent from traditional SAST tools.
    LLM-based analysis is best positioned as a triage aid or
    complementary signal rather than a standalone scanner.

  \item \textbf{Popularity does not predict better security.}
    Bandit and Semgrep show weak \emph{positive} correlations between
    star count and vulnerability density (Bandit: $r = 0.208$; Semgrep:
    $r = 0.134$), while the LLM shows a weak \emph{negative} correlation
    ($r = -0.139$). The inconsistent direction across tools suggests the
    relationship is fragile and tool-dependent. In all cases, popularity
    is not a reliable proxy for security quality.

  \item \textbf{Vulnerability density is independent of repository
    size and age.} Neither larger nor older repositories exhibit
    systematically different vulnerability densities, suggesting that
    security issues are introduced at relatively constant rates
    regardless of project maturity or scale.
\end{enumerate}

\subsection{Implications for Developers}

Our findings reinforce the importance of integrating SAST tools into
continuous integration pipelines. The presence of high-severity
findings in popular, heavily starred repositories---combined with the
weak positive correlation between popularity and finding
density---suggests that community trust (measured by stars) is not a
reliable proxy for security quality. We recommend that maintainers of
Python networking libraries adopt at least one SAST tool, enable
confidence-based filtering to suppress noise, and triage high-severity
findings as part of their release process.

\subsection{Validating the Multi-Tool Approach}

Our use of both Bandit and Semgrep directly addresses a gap identified
by prior literature. Alfadel et al.~\cite{Alfadel2023} concluded that
current security tools used by maintainers ``are not qualified to find
more complex and critical issues'' and called for robust tools that
combine multiple analysis techniques. By employing two SAST tools with
distinct rule profiles, our methodology provides broader detection
coverage than single-tool studies such as Ruohonen et
al.~\cite{Ruohonen2021}, pending the per-repository overlap analysis
that would quantify the marginal coverage each tool contributes.
Furthermore, while Alfadel et al.'s analysis relied on mining existing
vulnerability reports from third-party databases (Snyk.io), our
approach actively scans source code, enabling detection of potential
vulnerabilities that have not yet been reported or assigned CVE
identifiers.

\subsection{Ecosystem-Level Implications}

Alfadel et al.~\cite{Alfadel2023} recommended that the Python
ecosystem implement automated code vetting processes before packages
are published to PyPI. The \textsc{NetVulnPy} pipeline demonstrates
that such automated scanning is technically feasible at scale: our
pipeline processes thousands of repositories with configurable
parallelism and produces structured, queryable results. A similar
approach could be integrated into package registries or CI/CD
workflows to provide baseline security screening before code reaches
production.

\subsection{Threats to Validity}

\textbf{Internal validity.} SAST tools produce both false positives
and false negatives. To quantify false-positive rates, we validated a
stratified sample of 100~findings per tool using an independent LLM
verifier (Section~\ref{sec:results}). The estimated precision ranges
from 9.0\% (LLM) to 21.0\% (Semgrep), confirming that all reported
counts should be interpreted as \emph{potential} vulnerabilities.
However, our validation itself uses an LLM as the verifier rather than
human expert review, which introduces two potential biases: (1)~the
LLM verifier may share systematic blind spots with the LLM analyzer
(self-evaluation bias), potentially inflating or deflating precision
estimates for LLM findings; and (2)~the verifier may apply different
severity thresholds than a human security auditor. While we mitigated
bias~(1) by using a single verification model across all three tools
(ensuring consistent evaluation criteria), human validation of a
subset would further strengthen these estimates.

\textbf{External validity.} Our dataset is limited to repositories
discoverable via the GitHub Search API with specific
networking-related keywords. Keyword-based harvesting over-recovers web
frameworks and general-purpose libraries that are not strictly
networking tools, and we did not manually validate the domain
membership of the sampled repositories; the ``networking''
characterization should therefore be read as keyword-defined rather
than manually curated. The sample also excludes private repositories
and code hosted elsewhere, and sorting by stars introduces a
popularity bias.
%% TODO: consider manually coding a random sample of repositories to
%% estimate what fraction are genuinely networking-domain projects, and
%% report that domain-precision figure.

\textbf{Construct validity.} The mapping of Semgrep severities
(ERROR $\rightarrow$ HIGH, WARNING $\rightarrow$ MEDIUM,
INFO $\rightarrow$ LOW) is a simplification; Semgrep's severity
semantics differ from Bandit's, which limits the validity of direct
cross-tool comparisons of severity proportions.


%%═════════════════════════════════════════════════════════════════════
\section{Conclusion and Future Work}
\label{sec:conclusion}

We presented \textsc{NetVulnPy}, an automated pipeline for
large-scale vulnerability analysis of open-source Python networking
repositories. By applying Bandit, Semgrep, and an LLM-based analyzer
to 1{,}984 repositories, we identified 698{,}020 Bandit findings
(91{,}723 excluding the dominant low-severity \texttt{assert\_used}
rule), 5{,}614 Semgrep findings, and 82{,}741 LLM findings,
demonstrating that the three tools offer distinct rule profiles and
severity distributions. A precision validation on 100~stratified
samples per tool reveals that Semgrep achieves the highest precision
(21.0\%), followed by Bandit (10.0\%) and the LLM (9.0\%), with
HIGH-severity precision reaching 47.1\% for Semgrep. The LLM-based
analyzer contributes vulnerability categories requiring semantic
reasoning (SSRF, path traversal, context-aware credential detection)
but at the cost of taxonomic consistency and high false-positive rates.

By contrasting our first-party source-code findings against the
dependency-chain analysis of Alfadel et al.~\cite{Alfadel2023} and
the PyPI-wide scan of Ruohonen et al.~\cite{Ruohonen2021}, we
provide a more complete picture of the Python networking security
landscape.

Future work includes: (1)~recomputing vulnerability density with B101
excluded and on high-severity findings only, and applying
confidence-based filtering; (2)~complementing the LLM-based precision
validation with human expert review on a subset to calibrate the
automated verifier; (3)~a per-repository and per-file overlap analysis
to quantify Bandit/Semgrep/LLM complementarity; (4)~evaluating
additional LLM models and prompt strategies to improve taxonomic
consistency and reduce hallucinated findings; (5)~longitudinal
analysis tracking how vulnerability density evolves across repository
versions; (6)~extending the pipeline to support additional programming
languages.

%%─── Acknowledgments ─────────────────────────────────────────────────
\begin{acks}
This research was conducted in University of Passau. I am especially grateful to TA Talaya Farasat for her technical expertise and guidance throughout the research.

\end{acks}


%%─── References ──────────────────────────────────────────────────────
\bibliographystyle{ACM-Reference-Format}
\bibliography{references}

%%─── Appendix ────────────────────────────────────────────────────────
\appendix

\section{Pipeline Usage}
\label{app:usage}

The \textsc{NetVulnPy} pipeline is invoked as follows:

\begin{verbatim}
# Full pipeline with Bandit (default)
python main.py --max-repos 100 --verbose

# Full pipeline with Semgrep
python main.py --max-repos 100 --analyzer semgrep --verbose

# Full pipeline with LLM analyzer
python main.py --max-repos 100 --analyzer llm \
  --llm-api-key <key> --verbose

# Re-run analysis only (repos already downloaded)
python main.py --skip-harvest --skip-download --verbose

# Launch the interactive dashboard
streamlit run dashboard/app.py
\end{verbatim}


% \section{Database Schema}
% \label{app:schema}

% \begin{verbatim}
% CREATE TABLE repos (
%   full_name TEXT PRIMARY KEY,
%   stars INTEGER, forks INTEGER,
%   language TEXT, license TEXT,
%   created_at TEXT, updated_at TEXT,
%   size_kb INTEGER, ...
% );

% CREATE TABLE findings (
%   id INTEGER PRIMARY KEY AUTOINCREMENT,
%   repo TEXT REFERENCES repos(full_name),
%   filename TEXT, test_id TEXT,
%   test_name TEXT, issue_severity TEXT,
%   issue_confidence TEXT, issue_text TEXT,
%   line_number INTEGER, code TEXT
% );

% CREATE TABLE scan_summary (
%   full_name TEXT, py_files_found INTEGER,
%   total_issues INTEGER,
%   high INTEGER, medium INTEGER, low INTEGER,
%   exit_code INTEGER, analyzer TEXT,
%   status TEXT, loc INTEGER
% );
% \end{verbatim}

\end{document}